% Source: https://artofproblemsolving.com/community/c3h1617760_quick_2_second_problem
% Stable identity: aops-topic-1617760-q01
The characteristic equation $r^2=3r+2$, that is $r^2-3r-2=0$, has roots
\[
r_1=\dfrac{3+\sqrt{17}}{2},\qquad r_2=\dfrac{3-\sqrt{17}}{2}.
\]
They are distinct, so $a_n=Ar_1^n+Br_2^n$ for constants $A,B$ (any such expression satisfies the recurrence, and $A,B$ can be chosen to match $a_0,a_1$). From $a_0=2$: $A+B=2$. From $a_1=4$:
\[
Ar_1+Br_2=\dfrac{3(A+B)}{2}+\dfrac{\sqrt{17}(A-B)}{2}=3+\dfrac{\sqrt{17}(A-B)}{2}=4,
\]
so $A-B=\dfrac{2}{\sqrt{17}}$. Hence $A=1+\dfrac{1}{\sqrt{17}}$ and $B=1-\dfrac{1}{\sqrt{17}}$, and
\[
\boxed{a_n=\left(1+\dfrac{\sqrt{17}}{17}\right)\left(\dfrac{3+\sqrt{17}}{2}\right)^n+\left(1-\dfrac{\sqrt{17}}{17}\right)\left(\dfrac{3-\sqrt{17}}{2}\right)^n}.
\]
Check: $a_2=3\cdot4+2\cdot2=16$, and the formula gives $16$ as well.
